% ============================================================
%  DRAFT — Attention, Perception, & Psychophysics submission
%  "Pitch–Hue Correspondences: A Quantitative Meta-Analysis
%   of Three Centuries of Mapping Systems"
%  Status: Draft – July 2026
% ============================================================

\documentclass[12pt]{article}

%  -  Packages  -  -  -  -  -  -  -  -  -  -  -  -  -  -  - --
\usepackage[margin=1in]{geometry}
\usepackage{setspace}
\doublespacing
\usepackage{times}
\usepackage{natbib}
\setcitestyle{authoryear,open={(},close={)}}
\usepackage{amsmath}
\usepackage{booktabs}
\usepackage{array}
\usepackage{graphicx}
\usepackage{subcaption}
\usepackage{hyperref}
\usepackage{lineno}
\linenumbers

% Tell LaTeX where to look for figures
\graphicspath{
  {figures/}
  {figures/aesthetic/}
  {figures/aesthetic/families/}
  {figures/background/}
  {figures/litreview/}
}

% ============================================================
\begin{document}

%  -  Title page  -  -  -  -  -  -  -  -  -  -  -  -  -  -  - 
\begin{titlepage}
\begin{center}
\vspace*{2cm}
{\Large\bfseries Pitch--Hue Correspondences:\\[6pt]
A Quantitative Meta-Analysis of Three Centuries of Mapping Systems}\\[2em]

{\large Alex Noah Feldman and Ercan Engin Kuruoglu}\\[0.75em]

\textit{Draft -- July 2026}\\[0.5em]
\textit{Prepared for submission to \textbf{Attention, Perception, \& Psychophysics}}\\[3em]

\textbf{Running head:} PITCH--HUE CORRESPONDENCES META-ANALYSIS\\[3em]

% TODO(authors): fill in corresponding author, affiliation, and email
% before submission.
\textbf{Correspondence:}\\
Alex Noah Feldman\\
\texttt{[affiliation]}\\
\texttt{[email]}

\end{center}
\end{titlepage}

% ============================================================
\newpage
\begin{abstract}
Which color is the note C? Newton said red; Rimington said violet.
Three centuries of proposals have produced no consensus, and the
pattern of that disagreement is informative in its own right. We
present a quantitative meta-analysis of 70 pitch--hue mapping systems,
from Chinese Wuxing theory to contemporary digital applications, and
test four theoretical accounts of cross-modal correspondence:
structural (shared physics), semantic (shared language), statistical
(environmental co-occurrence), and hedonic (shared affective valence).
After standardizing all color specifications to CIELAB space and
applying the CIEDE2000 perceptual distance formula, global agreement
across the 25 retained systems is low (mean 0.22, per-pitch 0.20--0.30),
close to the simulated chance level of 0.21: the
historical record contains no universal pitch--hue mapping. Hue
assignments are nonetheless non-uniformly distributed (Rayleigh
$R = .66$, Monte Carlo $p = .003$), and the disagreement is patterned.
Organized into eight reasoning families classified from source texts,
the corpus shows a consistent within-family ordering:
Pedagogical systems agree most (0.348, 57\% above the pooled
baseline), followed by Synesthetic (0.276) and Spectral Analogy
(0.267) traditions, whereas Cosmological systems fall below the pooled
baseline---a descriptive ordering that small family sizes leave
statistically unresolved. Among the four accounts, the hedonic account
performs best: Procrustes alignment of cross-modal dimensions in
Evaluation-Potency-Activity affective space yields a disparity of
0.559 (permutation $p = .027$), with Evaluation as the primary
aligning dimension. The structural account earns partial credit
(semitone--hue distance correlation $r = .393$, permutation
$p = .002$), but its complementary-hue prediction fails.
Sentence-transformer probing recovers pitch--hue geometry under
timbral descriptors ($\rho = .368$, $p = .007$), consistent with a
semantics--affect entangled channel. Pitch--hue correspondences appear
organized by a shared affective geometry, scaffolded by structural and
statistical regularities; we derive from this a congruency framework
for neural research.
\end{abstract}

\noindent\textit{Keywords:} cross-modal correspondence, pitch--color
association, synesthesia, perceptual congruency, affective mediation

\vspace{1em}
\noindent\textit{Public Significance Statement:} Proposals linking
musical notes to colors have accumulated for three centuries without
agreement. Treating this historical record as data, we show that the
disagreement is structured: proposals agree within reasoning
traditions, and the traditions align best through the emotional
character shared by sounds and colors. The findings give educators,
accessibility designers, and neuroscientists a principled basis for
choosing sound--color mappings.

\newpage

% ============================================================
\section{Introduction}

Few questions in the science of perception have attracted such sustained
intellectual commitment - and produced such strong disagreement - 
as the relationship between musical pitch and color. Isaac Newton
divided the visible spectrum into seven color bands and aligned them
with the seven notes of the diatonic scale \citep{newton1704}.
Louis-Bertrand Castel built a keyboard instrument designed to project
colors alongside musical tones \citep{castel1725}. Alexander Scriabin
notated a full color-organ part in his orchestral score \textit{Prometheus:
The Poem of Fire} \citep{scriabin1911}. Olivier Messiaen described
vivid, involuntary color experiences triggered by specific chord
voicings \citep{messiaen1944}. Today, smartphone applications map sound
spectra to color in real time. Three centuries of effort by physicists,
composers, instrument builders, synesthetes, and software engineers
have all circled the same question - and arrived at irreconcilably
different answers. Scriabin and Rimsky-Korsakov, contemporaries who
both took the correspondence with utmost seriousness, agreed on little
beyond assigning warm colors to C \citep{galeyev2007}.

The disagreement is not for lack of musical reasoning - if anything,
each proposal is over-determined by it. Newton's seven spectral bands
were not chosen by eye: he deliberately partitioned the spectrum to
echo the seven unequal steps of the Dorian diatonic scale, importing
the ratios of the musical octave into the geometry of the prism
\citep{newton1704,gage1999}. Castel's \textit{clavecin oculaire}
literalized the analogy into a machine, pairing each semitone with a
pigment so that a harpsichordist could ``paint'' chords as they were
played \citep{castel1725,peacock1988}. A century later the color-organ
tradition - Rimington's electric \textit{Colour-Organ} foremost among
them - rebuilt the same dream with arc lamps and gelatin filters,
mapping the chromatic scale onto a calibrated hue wheel
\citep{rimington1912,peacock1988}. The composers who cared most
disagreed most instructively: Scriabin derived his \textit{luce} part
in \textit{Prometheus} from the circle of fifths, so that keys a fifth
apart received neighboring hues, whereas Rimsky-Korsakov reported
direct chromesthetic sensations that followed no such rule, and
Messiaen described seeing shifting complexes of color bound not to
single pitches but to specific chord voicings and modes
\citep{scriabin1911,messiaen1944,galeyev2007}. Older still, the
Chinese \textit{Wuxing} tradition linked the five pentatonic tones to
the five phases and their canonical colors, embedding pitch--color
correspondence inside a cosmological rather than perceptual logic
\citep{gage1999}. Each tradition is internally coherent; collectively
they are mutually contradictory.

This persistent disagreement is worth taking seriously as evidence.
Instead of asking which mapping is correct, we ask whether there is
measurable structure in three centuries of pitch--hue proposals, and if
so, which theoretical account best explains it. To do this we treat the
historical corpus itself as a dataset---something that, to our
knowledge, no prior survey has done quantitatively---and subject it to
statistical analysis.

The stakes extend well beyond historical curiosity. A principled,
data-driven account of pitch--hue correspondence would unlock concrete
applications across several fields. In \emph{music education}, color
coding is already widely used to teach pitch, notation, and harmony to
beginners (Boomwhackers, Figurenotes, and the colored-bell Suzuki and
Kod\'aly materials all assign colors to notes); knowing which mappings
are perceptually and affectively natural rather than arbitrary would
let educators design materials that learners find intuitive rather than
confusing. In \emph{accessibility and sensory substitution}, robust
audiovisual mappings underpin tools that render music visible for the
deaf and hard-of-hearing and sound audible-as-color for assistive
displays. In \emph{human--computer interaction and generative media},
real-time audio-to-color engines drive music visualizers, lighting
design, and the rapidly growing class of AI audiovisual generators,
all of which currently rely on ad hoc or designer-chosen palettes. And
in \emph{cognitive neuroscience}, a validated congruency criterion is a
prerequisite for any experiment that hopes to localize the neural
substrate of cross-modal binding - without one, ``congruent'' and
``incongruent'' stimuli are defined circularly by the very intuitions
under study. Resolving the structure of pitch--hue correspondence is
therefore not an end in itself but an enabling step for pedagogy,
assistive technology, creative tooling, and perceptual neuroscience
alike.

Cross-modal correspondences (CMCs) are systematic matching tendencies
between perceptual dimensions from different sensory domains, observed
reliably across individuals in the general population \citep{spence2011}.
Pitch--hue is among the most studied yet least settled CMC pairings.
Behavioral evidence is frustratingly weak: Ward, Huckstep, and
Tsakanikos (2006) found that participants associate higher pitches with
lighter, more yellow-shifted colors, but the effect is small, highly
variable, and entangled with brightness and saturation confounds
\citep{ward2006}. Rather than adding yet another explicit matching
experiment to an already crowded literature, we exploit a richer and
largely untapped empirical resource: three centuries of independent
pitch--hue proposals produced by agents as diverse as Pythagorean
cosmologists, chromesthetic synesthetes, music educators, and digital
application designers. Any signal these proposals carry should show up
in the structure of their agreement and disagreement, which we now have
the tools to measure.

Four accounts have been proposed to explain CMCs in general and
pitch--hue correspondences in particular \citep{spence2011,deroy2016}.
The \textbf{structural account} \citep{marks1978} holds that
correspondences arise from shared physical or mathematical properties:
both pitch and hue are circular dimensions with interval-ratio
relationships, and the pitch-class circle and the color wheel may share
the same underlying geometry. The \textbf{semantic account}
\citep{martino2001,walker2016} holds that shared linguistic labels
mediate correspondences: ``C major'' and ``red'' are associated because
Western cultural history has described both with overlapping emotional
vocabulary (``warm'', ``powerful'', ``grounded''). The
\textbf{statistical account} \citep{parise2014} holds that
correspondences arise from environmental co-occurrence
regularities - high-pitched bird calls in bright morning light,
low-frequency rumbles in dark overcast conditions - internalized through
Bayesian learning. The \textbf{hedonic account} \citep{palmer2013}
holds that stimuli evoking similar emotional responses are associated:
correspondence operates through a shared affective dimension onto which
both modalities project, rather than through direct physical or
linguistic links.

The present study tests all four accounts against a newly constructed
catalogue of 70 pitch--hue mapping systems, the largest assembled to
date. The approach differs from prior work in three ways. First, instead
of adopting any single reference map as ground truth, we quantify the
\emph{geometry of agreement} across systems, asking whether musical
interval structure is preserved in color distance structure; this
reference-free framing sidesteps the circularity of earlier comparisons.
Second, we organize systems into eight theoretically motivated
\emph{reasoning families}---classified from source texts before any
agreement is computed---and test whether a shared reasoning tradition
predicts shared color assignment. This lets us test whether the global
disagreement seen in prior surveys conceals coherent within-tradition
structure---a form of meta-analytic heterogeneity in which pooling
masks subgroup agreement. Third, we probe the semantic
account directly using pretrained sentence-transformer models
\citep{wang2020,song2020,nussbaum2024} trained on large
natural-language corpora, turning the linguistic mediation hypothesis
into a directly testable claim.

A note on the kind of paper this is: the contribution is meta-analytic
and methodological. We present no new behavioral data. What we offer
is a curated catalogue, a reference-free agreement framework, and a
head-to-head comparison of theoretical accounts---quantitative targets
for the behavioral and neural experiments that should follow.

% ============================================================
\section{Materials: Constructing the Historical Catalogue}

\subsection{Source Identification and Coverage}
Primary sources were identified through three channels. First, four
academic surveys of the sound--color relationship
\citep{jewanski2009,gage1999,peacock1988,wells1980} were reviewed
systematically to extract every referenced pitch--color mapping system.
Second, targeted searches in music history archives, color science
literature, and instrument design records extended coverage to include
practitioner systems - color organs, music education curricula, digital
applications - underrepresented in academic surveys. Third, a
bibliographic snowball procedure followed citations until no new systems
emerged. The procedure yielded 70 candidate systems spanning from
Chinese Wuxing theory (4th century BCE) to contemporary smartphone
color--music applications.

\subsection{Color Standardization}
Historical color specifications are heterogeneous: Newton specifies
spectral band positions; Scriabin specifies color names in Russian;
Rimington specifies 19th-century pigment trade names; digital
applications specify hexadecimal RGB codes; synesthetes report colors
in free text. Each specification was classified into one of four source
types: (1)~spectral wavelength, (2)~named pigment or dye,
(3)~natural-language description, or (4)~modern color code (RGB/hex),
and converted to CIELAB coordinates via a principled pipeline. Spectral
wavelengths were converted to CIEXYZ using the CIE 1931 standard
observer, then to CIELAB. Named pigments and dyes were matched to
Munsell \citep{newhall1943} or Natural Color System
\citep{hard1981} references using published pigment databases.
Natural-language descriptions were resolved by averaging the CIELAB
coordinates of all XKCD Color Survey \citep{munroe2010xkcd} colors
judged semantically compatible with the description. Modern color codes
were converted via the standard sRGB-to-CIELAB transformation with D65
illuminant.

\subsection{Inclusion Criteria}
Each candidate system was evaluated against four criteria:
(1)~\textit{Discrete pitch-class-to-color mapping}: the system must
assign colors to discrete pitch classes. Systems mapping intervals,
modes, timbres, keys that cannot be normalized to the 12 pitch
classes, or continuous frequency were excluded, as were systems
specifying only lightness or brightness without hue. Partial coverage
(pentatonic, diatonic, or heptatonic systems) was retained; per-pitch
statistics are computed over the systems that specify that pitch
class, and interpolation of missing pitch classes was rejected because
any interpolation scheme imposes structural assumptions that the
analysis is designed to test.
(2)~\textit{Convertible color specifications}: all assignments must
be convertible to CIELAB coordinates.
(3)~\textit{Traceable provenance}: each system must be attributable to
a published source with sufficient detail to verify color assignments.
(4)~\textit{Independence}: directly derived or copied systems were
merged with their originals.

After applying these criteria, 25 of the candidate systems identified
through the survey channels were retained for the primary quantitative
analyses; 14 of the retained systems specify all 12 chromatic pitch
classes. The remaining catalogued systems are documented but excluded
from agreement computations. Supplementary Table~S1 lists all 70
systems with year, source, family, specification type, retention
status, and exclusion reason.
Six further systems meeting the same criteria entered the catalogue
through the snowball procedure after the primary analyses were run,
bringing the total of retained systems in the final catalogue to 31.
For two of them, Castel's \textit{clavecin oculaire} (1725) and
Rimington's color-organ assignments (1895)---both discussed, but not
fully tabulated, in prior surveys \citep{peacock1988,gage1999}---%
complete 12-tone color specifications were assembled from the primary
sources. The remaining four are Field's chromatic equivalents
\citep{field1835}, Wells' proportional frequency system
\citep{wells1980}, Pridmore's psychophysical transposition
\citep{pridmore1992}, and a contemporary smartphone Soundmap
application. Unless noted otherwise, results below are reported for
the 25-system primary set; a sensitivity analysis on all 31 retained
systems (flat agreement 0.237, Rayleigh $R = .60$, semitone--hue
correlation $r = .49$) leaves every conclusion unchanged.

\subsection{Reasoning-Family Classification}
The 70 systems were classified into eight reasoning families (F1--F8)
based on the design rationale of each system \emph{before} agreement
scores were computed, to avoid circular analysis. Two independent
coders assigned each system to a family from the original source text;
disagreements ($\approx$12\% of cases) were resolved through discussion.
The eight families are: F1~Spectral Analogy (22 systems), F2~Circle of
Fifths (4), F3~Cosmological (14), F4~Emotion-Mediated (3),
F5~Synesthetic (7), F6~Esoteric (2), F7~Psychophysical (8), and
F8~Pedagogical (4); the remaining six catalogue entries are meta-level
sources (survey compilations and documented rejections of fixed
mappings) that name no single design rationale and were left
unclassified. Twelve systems invoke a secondary rationale in addition
to their primary one; the priority rule described below assigns each
to a single primary family.

The eight categories follow directly from the four theoretical accounts
the study sets out to test, refined to the granularity actually present
in the historical record. The \emph{structural} account corresponds to
F1~(Spectral Analogy, Newton's prism metaphor) and F2~(Circle of
Fifths, harmonic interval geometry); the \emph{hedonic} account to
F4~(Emotion-Mediated) and, indirectly, F8~(Pedagogical), where
designers optimize for what feels natural to a learner; the
\emph{semantic} and culturally transmitted accounts to
F3~(Cosmological, pitch--color links inherited from metaphysical
systems) and F6~(Esoteric, Theosophical and occult schemes).
F5~(Synesthetic) isolates first-person chromesthetic report, and
F7~(Psychophysical) isolates systems derived from explicit perceptual
measurement. Each family is defined by a single text-based inclusion
criterion---the \emph{stated} design rationale of the system, not its
resulting colors (Table~\ref{tab:family_definitions}). Coders applied
the criteria in fixed order F1 through F8, assigning each system to the
first family whose criterion it met; this priority rule resolves the
small number of systems invoking more than one rationale (a color organ
justified on both spectral and pedagogical grounds is coded F1, its
primary stated basis). Eight is the smallest number of families that
keeps the four accounts distinguishable while separating the two
recurring expressions of structural and hedonic reasoning that the
record contains; finer splits produced singleton families, and coarser
ones merged rationales the sources treat as distinct. Retaining the
eight-way scheme rather than collapsing to the four accounts lets a
single account succeed in one guise and fail in another (structural
reasoning expressed spectrally versus harmonically, for example)---a
distinction that proves diagnostic.

\begin{table}[htbp]
\centering
\caption{Operational definitions of the eight reasoning families. Each
system was assigned to the first family (in order F1--F8) whose stated
design rationale it satisfied. Counts are catalogue-level (of the 70
candidate systems); the subset retained for agreement analysis is
smaller (see Section~\ref{sec:results} and
Table~\ref{tab:family_agreement}).}
\label{tab:family_definitions}
\renewcommand{\arraystretch}{1.5}
\setlength{\tabcolsep}{6pt}
\begin{tabular}{@{}p{0.8cm}p{3.0cm}p{8.2cm}p{2.9cm}@{}}
\toprule
\textbf{ID} & \textbf{Family} & \textbf{Defining criterion (stated
  rationale)} & \textbf{Representative system} \\
\midrule
F1 & Spectral Analogy & Color assigned by analogy to the optical
  spectrum or the rainbow (Newton's prism metaphor) & Newton (1704);
  Castel (1725) \\
F2 & Circle of Fifths & Color assigned to preserve harmonic interval
  geometry, neighboring keys receiving neighboring hues & Scriabin
  (1911) \\
F3 & Cosmological & Color inherited from a metaphysical or
  cosmological system (phases, planets, elements) & Chinese
  \textit{Wuxing} \\
F4 & Emotion-Mediated & Color chosen to match the explicitly named
  emotional character of the pitch or key & Mood-based key--color
  charts \\
F5 & Synesthetic & Color reported as a first-person, involuntary
  chromesthetic experience & Rimsky-Korsakov; Messiaen (1944) \\
F6 & Esoteric & Color derived from Theosophical, occult, or spiritual
  doctrine & Theosophical scales \\
F7 & Psychophysical & Color derived from explicit perceptual or
  psychophysical measurement & Pridmore (1992) \\
F8 & Pedagogical & Color chosen to be intuitive for teaching pitch or
  notation to learners & Boomwhackers; colored-bell methods \\
\bottomrule
\end{tabular}
\end{table}

\subsection{Color Distance Metric}
\label{sec:metric}

All inter-color distance computations used the CIEDE2000 formula
\citep{sharma2005}:
\begin{equation}
\Delta E_{00} = \sqrt{
  \left(\frac{\Delta L'}{k_L S_L}\right)^2 +
  \left(\frac{\Delta C'}{k_C S_C}\right)^2 +
  \left(\frac{\Delta H'}{k_H S_H}\right)^2 +
  R_T \frac{\Delta C'}{k_C S_C} \frac{\Delta H'}{k_H S_H}
}
\label{eq:ciede2000}
\end{equation}
where $\Delta L'$, $\Delta C'$, $\Delta H'$ are lightness, chroma, and
hue differences; $S_L$, $S_C$, $S_H$ are empirically calibrated
weighting functions; $k_L = k_C = k_H = 1$; and $R_T$ corrects
perceptual non-uniformities in the blue-violet region. CIEDE2000 was
preferred over raw CIELAB Euclidean distance to ensure that differences
in the blue--green region are not artifactually underweighted. An
\textbf{agreement score} was derived as the inverse of normalized mean
pairwise CIEDE2000 distance:
\begin{equation}
  \text{agreement} = \frac{1}{1 + \overline{\Delta E_{00}} / 12}
\label{eq:agreement}
\end{equation}
The denominator of 12 was calibrated so that a score of 0.60
corresponds to a mean CIEDE2000 distance of $\approx$8 units, a
conventional benchmark above which color pairs are judged clearly
distinct under standard viewing conditions \citep{sharma2005}; we
treat 0.60 as an explicitly chosen benchmark for ``strong consensus''
rather than as a perceptual constant. A chance baseline for the
agreement score was estimated by passing simulated systems of 12
colors, drawn uniformly from the sRGB gamut, through the identical
pipeline (2,000 simulated ensembles per corpus size). For ensembles of
24 random systems the simulated chance agreement is 0.212 (95\%
interval $[0.196, 0.231]$); for ensembles of 15 it is 0.212
$[0.192, 0.240]$, with wider intervals at smaller ensemble sizes
(full grid in Supplementary Table~S2). Observed scores are compared
against this simulated distribution wherever ``chance'' is invoked
below.

One implementation detail matters for the hue-angle analyses reported
below. Consensus hue angles were computed by averaging the CIELAB
$a^*$/$b^*$ coordinates across systems and taking the hue angle of the
mean vector. This is equivalent to a chroma-weighted circular mean: it
prevents instability near the achromatic axis by construction, since
near-achromatic assignments (e.g., Scriabin's ``steely'' grays)
contribute vectors of small magnitude and thus receive
correspondingly little weight.

% ============================================================
\section{Analyses and Results}
\label{sec:results}

The four theoretical accounts are taken up in turn, each method stated
immediately before the result it yields. The scoreboard, in brief: the
hedonic account performs best (Procrustes disparity 0.559 in affective
space); the structural account earns partial credit (interval ordering,
but not specific assignment); the statistical account passes its
distinctive variance prediction; and the semantic account contributes a
channel that cannot be fully separated from affect. Four claims were
designated as primary---family-level heterogeneity, the EPA Procrustes
alignment, the interval--variance correlation, and the timbral
embedding alignment---and multiplicity across the full test battery is
addressed in the Discussion and Supplementary Table~S7. All agreement
computations rest on the CIEDE2000 distance and agreement score defined
in Equations~\ref{eq:ciede2000} and~\ref{eq:agreement}.

\subsection{Three Centuries, Seventy Systems}

The 70 catalogued systems span 23 centuries of intellectual history;
the retained corpus is concentrated in the three centuries from Diez
(1723) and Castel (1725) to the present---the ``three centuries'' of
our title---with three ancient cosmological systems (Chinese
\textit{Wuxing}, Ancient Persian, Indian svara) as the only earlier
retained entries. The
temporal distribution is itself revealing. Figure~\ref{fig:timeline}
shows that proposal density peaks sharply in the 19th and early 20th
centuries - the era of color organs and synesthetic aesthetics - and
surges again in the digital age, with a pronounced trough from
approximately 1930 to 1980 when serialist and atonal movements
displaced color from the mainstream musical imagination. Newton's
shadow is long: the Spectral Analogy family (F1) dominates the corpus
(Figure~\ref{fig:timeline}, right), with his prism metaphor driving
more than a third of all proposals across three centuries.

Of the 70 candidate systems, 31 met all four inclusion criteria (25
in the primary set); the dominant sources of exclusion were mappings
defined over intervals, modes, or timbres rather than pitch classes,
and specifications restricted to lightness without hue. That Newton's
own 7-note scheme left the five chromatic pitch classes unassigned is
a choice that cascaded through centuries of followers.

As a check on the hand-coded taxonomy, we clustered the systems of
the merged family corpus (Section~\ref{sec:results}) by their pairwise
CIEDE2000 distance profiles using Ward-linkage hierarchical clustering
\citep{ward1963}. The resulting partition recovers the hand-coded
families only weakly (adjusted Rand index $= .23$ at $k = 5$, $.12$ at
$k = 8$). This is informative rather than disqualifying: the families
were coded from \emph{stated design rationale}, not from the colors
themselves, and the low Rand index confirms that the two codings are
not redundant---systems within a reasoning tradition need not produce
similar palettes, which is precisely the question the family-level
agreement analysis quantifies.

\begin{figure}[htbp]
\centering
\begin{subfigure}[t]{0.62\textwidth}
  \centering
  \includegraphics[width=\textwidth]{P1_01_timeline.png}
  \caption{Historical timeline of pitch--hue mapping proposals (Newton
  1704 to present). Proposal density peaks in the 19th--early 20th
  century and again in the digital age.}
\end{subfigure}
\hfill
\begin{subfigure}[t]{0.34\textwidth}
  \centering
  \includegraphics[width=\textwidth]{P1_01_family_distribution.png}
  \caption{Distribution of the 70 systems across the eight reasoning
  families. F1 (Spectral Analogy) is the largest.}
\end{subfigure}
\caption{Overview of the historical catalogue. Left: timeline.
Right: family distribution.}
\label{fig:timeline}
\end{figure}

\subsection{Global Consensus Is Weak but Non-Random}
\label{sec:global}

The global picture is clear: there is no universal pitch--hue mapping
in this corpus. CIEDE2000 agreement across all 25 retained systems
ranged from 0.20 to 0.30 across the 12 pitch classes (mean 0.221),
well below the 0.60 benchmark for strong consensus and close to the
simulated chance level of 0.212 for random color systems
(Section~\ref{sec:metric}); only G$\sharp$ (0.30) exceeds the upper
bound of the simulated chance interval. Every pitch class
spans a broad arc of hue space when all traditions are pooled
(Figure~\ref{fig:consensus}, left), so the data do not support any claim
that a particular system captures the ``true'' pitch--hue
correspondence.

The assignments are nonetheless far from random. The Rayleigh test,
applied to the 12 consensus hue angles (chroma-weighted circular means
across the 25 retained systems), returned $R = 0.66$; a Monte Carlo
test against 12 hues drawn uniformly on the circle gives $p = .003$.
The mean resultant direction is $24^\circ$, in the orange--red region:
a clear warm-color bias.
The pattern of variance is also informative (Figure~\ref{fig:consensus},
right): pitch classes C and G, the tonal anchors of Western harmony,
show the lowest inter-system spread, while the chromatic alterations
(C$\sharp$, D$\sharp$, F$\sharp$, G$\sharp$, A$\sharp$) show the
highest---a trace of Western tonal hierarchy in a corpus dominated by
Western sources.

\begin{figure}[htbp]
\centering
\begin{subfigure}[t]{0.46\textwidth}
  \centering
  \includegraphics[width=\textwidth]{P1_02_consensus_wheel.png}
  \caption{Global consensus wheel. Each wedge is positioned at the
  chroma-weighted circular-mean hue for one pitch class, and wedge
  opacity encodes the agreement score. The uniformly low opacities and
  the wide dispersion of mean hues reflect the absence of a reliable
  global consensus.}
\end{subfigure}
\hfill
\begin{subfigure}[t]{0.50\textwidth}
  \centering
  \includegraphics[width=\textwidth]{P1_02_delta_e_bars.png}
  \caption{Per-pitch-class CIEDE2000 standard deviation across all 25
  retained systems. C and G show the lowest variance (tonal anchors);
  chromatic alterations show the highest variance.}
\end{subfigure}
\caption{Global consensus analysis. Pooled agreement is low, but hue
assignments are non-uniformly distributed (Rayleigh $R = 0.66$,
Monte Carlo $p = .003$), indicating that the mapping is structured but
tradition-specific rather than universal.}
\label{fig:consensus}
\end{figure}

\subsection{The Eight Families: Structure Within Traditions}

The family-level analysis is where the paper's central descriptive
claim is established. This analysis uses a \emph{merged family corpus}:
the retained catalogue systems plus independently digitized versions of
several historical systems from our initial source collection, giving
39 systems across the five families with at least two members
(Table~\ref{tab:family_agreement}). Three systems in F1 appear twice
(Castel, Field, and Rimington, each digitized from two editions);
deduplicating them changes the F1 score by less than 0.001, so the
merged counts are reported as analyzed. Breaking the corpus into
reasoning families reveals a consistent ordering that the flat
analysis hides. The Pedagogical family (F8) leads
at 0.348, 57\% above the global consensus. Synesthetic systems (F5)
follow at 0.276, and Spectral Analogy (F1) at 0.267. The Cosmological
family (F3) goes the other way: it falls \emph{below} the flat baseline
at 0.187 vs.\ 0.221.

To quantify the uncertainty of these small-sample scores, each family
estimate is accompanied by a bootstrap confidence interval (systems
resampled within family, 10,000 draws, resamples constrained to
contain at least two distinct systems), and the within-family
advantage was tested by a family-label permutation test in which the
family labels were shuffled across the merged corpus and the family
agreement scores recomputed (10,000 permutations). The outcome is
sobering and worth stating plainly: no single family's deviation from
the pooled baseline is individually significant (F1 $p = .10$, F8
$p = .33$, F5 $p = .40$; F3's below-baseline score has $p = .91$
against the greater-than null, i.e., it sits in the lower tail as
predicted), and a global heterogeneity permutation test---the
label-shuffle null for the variance of family means---is likewise
non-significant ($p = .81$). With 3--5 systems in four of the five
families, the family ordering reported here is a descriptive,
effect-size-level pattern whose confirmation requires an expanded
corpus. Interval estimates and permutation $p$-values are reported in
Supplementary Table~S2.

\begin{table}[htbp]
\centering
\caption{Within-family versus flat-catalogue CIEDE2000 agreement
scores in the merged family corpus; $N$ is the number of systems
contributing to each family's score. Within-family agreement is
reported for families with $N \ge 2$ systems; F2~(Circle of Fifths),
F4~(Emotion-Mediated), and F7~(Psychophysical) each contributed a
single system and are omitted. F8, F5, and F1 exceed the flat
baseline; F6 and F3 fall below it. The flat baseline is computed on
the 25-system primary set. Bootstrap intervals and permutation
$p$-values are given in Supplementary Table~S2.}
\label{tab:family_agreement}
\begin{tabular}{@{}llccc@{}}
\toprule
\textbf{Family} & \textbf{Tradition} & \textbf{$N$} &
  \textbf{Agreement} & \textbf{$\Delta$ vs.\ flat} \\
\midrule
F8 & Pedagogical      & 3  & \textbf{0.348} & $+0.127$ \\
F5 & Synesthetic      & 5  & 0.276          & $+0.055$ \\
F1 & Spectral Analogy & 24 & 0.267          & $+0.046$ \\
F6 & Esoteric         & 3  & 0.201          & $-0.020$ \\
F3 & Cosmological     & 4  & 0.187          & $-0.034$ \\
\midrule
\multicolumn{2}{l}{\textit{Flat (all retained systems)}} & 25 & 0.221
  & -- \\
\bottomrule
\end{tabular}
\end{table}

Figure~\ref{fig:families} illustrates these results visually.
The bar chart (left) plots within-family agreement against the flat
baseline; the wheel panel (right) shows the mean hue assignment for
each pitch class within each tradition - confirming that these are not
minor quantitative differences but categorically distinct
color philosophies.

\begin{figure}[htbp]
\centering
\begin{subfigure}[t]{0.45\textwidth}
  \centering
  \includegraphics[width=\textwidth]{P1_02v2_family_bars.png}
  \caption{Within-family agreement scores compared against the flat
  baseline (dashed line). F8 and F5 substantially exceed global
  agreement; F3 falls below it.}
\end{subfigure}
\hfill
\begin{subfigure}[t]{0.52\textwidth}
  \centering
  \includegraphics[width=\textwidth]{P1_02v2_family_wheels_panel.png}
  \caption{Per-family consensus wheels. Point size is proportional to
  within-family agreement. The traditions assign color through
  fundamentally different color philosophies.}
\end{subfigure}
\caption{Family-level analysis. Within-family structure substantially
exceeds global agreement for Pedagogical and Synesthetic traditions,
while Cosmological systems fall below the flat baseline.}
\label{fig:families}
\end{figure}

The Cosmological family's result deserves emphasis: it is the most
important negative finding in the dataset. Centuries of cosmological
theorizing produce \emph{more} disagreement than pooling all retained
systems together. The four analyzed F3 systems span genuinely distinct
cosmological traditions---the Chinese \textit{Wuxing} five-phase
system, the Ancient Persian system, the Indian svara associations, and
Diez's planetary note--color scheme (1723)---so the low within-family
agreement reflects divergence between independent metaphysical
frameworks, not noise within one. This argues against any account that
treats cosmological reasoning as a primary driver of pitch--hue
correspondences.

The Pedagogical family points the other way, but two alternative
explanations must be addressed before reading its convergence as
affective. The first is lineage: commercial teaching materials copy one
another, so inclusion criterion~4 (independence) was applied strictly
to F8, merging derivative products with their sources. The three
analyzed pedagogical systems---Sudre's Solresol color code (France,
1862), Boomwhackers (United States, 1995), and Figurenotes (Finland,
1996)---arose in different countries, markets, and eras with no
documented lineage among them.
The second is structure: educators might simply reuse rainbow order,
which would make the convergence spectral rather than hedonic. We
therefore compare the F8 consensus mapping against the canonical
spectral-order mapping in Supplementary Analysis~S3; the two are close
(mean angular distance $35^\circ$ after optimal rotation), so the
pedagogical advantage should be read as convergence
on a spectrally anchored, affectively comfortable mapping rather than
as purely affective agreement. With these caveats stated, the
convergence remains striking: educators across eras, optimizing for
what feels natural to untrained learners, agree more than any other
tradition. Figure~\ref{fig:families_panel} shows the complete mapping
panel organized by family.

\begin{figure}[htbp]
\centering
\includegraphics[width=\textwidth]{P1_families_panel.png}
\caption{Within-family mapping panel. Each row shows one historical
system's color assignments for all 12 pitch classes (C to B). Families
are grouped by reasoning tradition (F1--F8). The Spectral Analogy
family (F1) is dominated by rainbow order; the Synesthetic family (F5)
tends toward warm hues at low pitches and cool hues at high pitches;
the Pedagogical family (F8) shows the tightest within-family
clustering.}
\label{fig:families_panel}
\end{figure}

The traditions also differ in \emph{how} they encode harmonic
relationships, not just in which colors they choose
(Figure~\ref{fig:family_ranking}). The Spectral Analogy and Synesthetic
families embed substantially stronger interval-to-hue structure
($\rho \approx 0.6$) than the Esoteric family, and all three assign more
similar hues to consonant than to dissonant intervals - an early hint of
the structural signal we quantify next.

\begin{figure}[htbp]
\centering
\includegraphics[width=0.82\textwidth]{P1_03v2_family_ranking.png}
\caption{Per-family harmonic structure analysis, shown for the three
families (Esoteric F6, Spectral Analogy F1, Synesthetic F5) for which
per-interval hue statistics could be computed on the retained systems.
Left: harmonic
structure score (mean hue distance for dissonant intervals minus
consonant intervals); all three families show a negative score,
indicating that consonant intervals are assigned more similar hues than
dissonant ones. Right: Spearman $\rho$ between semitone distance and
mean hue distance; Spectral Analogy and Synesthetic families show
substantially stronger structure ($\rho \approx 0.6$) than the Esoteric
family.}
\label{fig:family_ranking}
\end{figure}

\subsection{Testing the Structural Account: Does Interval Geometry
Predict Hue Geometry?}

\textbf{Method.} We tested six interval-to-hue hypotheses, ordered from
the strongest and most specific claim to the weakest
(Table~\ref{tab:structural_h}). The interval--distance correlation (H2)
is computed over the $\binom{12}{2} = 66$ pitch-class pairs of the
pooled consensus. Two properties of this design constrain inference:
circular pitch-class distance collapses the eleven interval sizes to
six distinct values (1--6 semitones), and pairwise distances are not
independent. Statistical significance was therefore assessed by a
pitch-label permutation test in addition to parametric inference: the
12 pitch-class labels were randomly reassigned across color
assignments 10,000 times to generate a null distribution of
correlation values.
The point-prediction hypotheses (H1, H3--H6) are evaluated
descriptively, by comparing the observed mean hue distance for the
relevant interval against the predicted angular range. We then
evaluated six principled
structural mapping theories (Table~\ref{tab:structural_theories}), each
a complete first-principles recipe for assigning hues to pitches, on
two complementary statistics. The first is a reference-free
\emph{geometry alignment} score: the Spearman correlation between the
$\binom{12}{2}$ pairwise circular pitch-step distances and the
corresponding pairwise hue angular distances of the theory's own
predicted mapping,
\begin{equation}
  \rho_{\text{geom}} = \operatorname{Spearman}\!\left(
  \left\{ d_{\text{pitch}}(i,j) \right\}_{i<j},
  \left\{ \left| \hat{h}_i - \hat{h}_j \right|_{\text{circ}}
  \right\}_{i<j} \right)
\label{eq:geometry}
\end{equation}
where $\hat{h}$ are the predicted hue angles and
$|\cdot|_{\text{circ}}$ denotes circular (wrap-around) angular
distance. This statistic asks whether the theory realizes an
interval-to-arc isomorphism at all, without privileging any historical
map as ground truth; its chance level was estimated from 10,000 random
hue assignments (null mean $\rho_{\text{geom}} \approx 0$, 95th
percentile $0.23$). The second statistic is the \emph{consensus
angular error}: the mean circular distance between the theory's
predicted hue and the observed pooled-consensus hue across the 12
pitch classes. Because two hue angles drawn independently at random
differ by $90^\circ$ on average, a theory unrelated to the data scores
$\approx 90^\circ$ on this measure (the \emph{random-hue null}); exact
Monte Carlo $p$-values are reported against that null. For calibration
we also report two external benchmarks: a perfect spectral-order
mapping ($\rho_{\text{geom}} = 1.00$) and the circle-of-fifths mapping
($\rho_{\text{geom}} = -.19$). That the spectral mapping attains the
maximal geometry score is itself a finding: the pooled historical
consensus preserves spectral hue order almost exactly, so the spectral
mapping functions as the effective ceiling of the scale.

\begin{table}[htbp]
\centering
\caption{Six structural hypotheses linking musical intervals to
color-wheel geometry, with outcomes. H1--H6 span from the strongest
claim (perfect fifth maps to complementary hue at 180$^\circ$) to the
weakest (monotone positive correlation between interval size and hue
distance). Verdicts for the point predictions (H1, H3--H6) are
descriptive comparisons of the observed mean hue distance
(Table~\ref{tab:harmonic}) against the predicted range; only H2
receives formal (permutation) inference. Note that circular
pitch-class distance makes the perfect fourth and perfect fifth
equivalent (both distance 5), so H1 and H5 address the same observed
value.}
\label{tab:structural_h}
\begin{tabular}{@{}p{0.5cm}p{2.5cm}p{1.6cm}p{5.2cm}p{2.8cm}@{}}
\toprule
\textbf{H} & \textbf{Interval} & \textbf{Semitones} &
  \textbf{Prediction} & \textbf{Result} \\
\midrule
H1 & Perfect fifth  & 7     & Complementary hues ($\sim$180$^\circ$):
  most consonant interval $\leftrightarrow$ maximal color contrast
  & Fail (80.7$^\circ$) \\
H2 & All intervals  & 1--11 & Semitone distance correlates positively
  with hue distance (interval-to-arc isomorphism)
  & \textbf{Pass} ($r = .393$, perm.\ $p = .002$) \\
H3 & Tritone        & 6     & Complementary hues ($\sim$180$^\circ$):
  most dissonant interval $\leftrightarrow$ maximal color contrast
  & Fail (85.7$^\circ$) \\
H4 & Major third    & 4     & Triadic hue separation (90--150$^\circ$,
  centered on 120$^\circ$): three major thirds tile the octave,
  mirroring three-fold color-wheel symmetry & Fail (70.5$^\circ$) \\
H5 & Perfect fourth & 5     & Intermediate hue distance
  (70--140$^\circ$): between analogous and complementary
  & \textbf{Pass} (80.7$^\circ$) \\
H6 & Minor second   & 1     & Adjacent hues ($<$60$^\circ$): chromatic
  neighbors $\leftrightarrow$ analogous, closely spaced colors
  & \textbf{Pass} (36.9$^\circ$) \\
\bottomrule
\end{tabular}
\end{table}

\begin{table}[htbp]
\centering
\caption{Six principled structural theories, evaluated on geometry
alignment (Equation~\ref{eq:geometry}) and on angular error against
the pooled consensus.}
\label{tab:structural_theories}
\begin{tabular}{@{}p{0.8cm}p{3.2cm}p{7.8cm}@{}}
\toprule
\textbf{Abbr.} & \textbf{Name} & \textbf{Core principle} \\
\midrule
IRP & Interval-Ratio Preservation    & Maps 12 pitch classes onto the color
  wheel so that semitone ratios mirror hue-angle ratios, optimizing rotation
  to maximize the consonance--color-distance correlation. \\[4pt]
FRI & Frequency-Ratio Isomorphism    & Maps C4--B4 on a log-frequency scale
  to the visible spectrum (380--780~nm), revising Newton's linear mapping for
  the logarithmic nature of pitch perception. \\[4pt]
HEM & Harmonic Entropy Matching      & Pairs pitches with low harmonic entropy
  (consonant, simple ratios) to focal colors \citep{berlin1969} on the
  grounds that shared perceptual salience mechanisms link consonance and
  color centrality. \\[4pt]
MPA & Multidimensional Perceptual Alignment & Constrains pitch height
  (mel scale) to lightness and spectral centroid to color temperature;
  hue is solved as the residual dimension. \\[4pt]
PDM & Perceptual Distinctiveness     & Maps tonal hierarchy positions
  (probe-tone ratings; \citealp{krumhansl1990}) to color naming
  consensus, pairing perceptually central pitches with the most-named
  colors. \\[4pt]
SSC & Spectral Shape Correspondence  & Aligns RDMs of piano harmonic spectra
  with hue spectral profiles, testing whether the physical frequency content
  of a tone predicts the hue whose spectral reflectance it most resembles. \\
\bottomrule
\end{tabular}
\end{table}

\textbf{Results.} The structural account earns a qualified pass. Of
the six interval hypotheses, the weakest (H2, monotone ordering), the
most local (H6, adjacent hues for semitone neighbors), and the
broad-band intermediate prediction (H5) pass; every point prediction
targeting a \emph{specific} angular value---complementary or
triadic---fails (Table~\ref{tab:structural_h}). The positive
correlation between
semitone distance and mean hue distance ($r = .393$, parametric
$p = .001$, pitch-label permutation $p = .002$) is
clear: across all retained systems, larger musical intervals
consistently map to larger color distances
(Figure~\ref{fig:structural}, left). Table~\ref{tab:harmonic} gives
the full breakdown by interval.

\begin{table}[htbp]
\centering
\caption{Mean hue distance (degrees on the color wheel) by musical
interval, pooled across all 25 retained systems. Circular pitch-class
distance makes the perfect fourth and perfect fifth equivalent (both
distance 5); their rows are identical by construction.}
\label{tab:harmonic}
\begin{tabular}{@{}lccc@{}}
\toprule
\textbf{Interval} & \textbf{Semitones} &
  \textbf{Mean hue distance ($^\circ$)} & \textbf{SD ($^\circ$)} \\
\midrule
Minor second   & 1 & 36.9 & 44.6 \\
Major second   & 2 & 53.3 & 44.3 \\
Minor third    & 3 & 59.1 & 38.3 \\
Major third    & 4 & 70.5 & 34.9 \\
Perfect fourth & 5 & 80.7 & 26.1 \\
Tritone        & 6 & 85.7 & 22.0 \\
Perfect fifth  & 7 & 80.7 & 26.1 \\
\bottomrule
\end{tabular}
\end{table}

The account's strongest prediction, however, fails. The claim that
harmonically complementary intervals should map to visually
complementary hues ($\approx$180$^\circ$ apart) is rejected
(Figure~\ref{fig:structural}, right): perfect fifths averaged only
80.7$^\circ$ of hue separation, less than half the predicted value. The
interval--distance relationship saturates near 85$^\circ$ and never
approaches the complementary region. Interval structure can account for
the ordinal pattern, then, but not for the specific mapping.

\begin{figure}[htbp]
\centering
\begin{subfigure}[t]{0.48\textwidth}
  \centering
  \includegraphics[width=\textwidth]{P1_03_semitone_hue_scatter.png}
  \caption{Semitone distance vs.\ mean hue distance across all
  historical systems ($r = .393$, $p = .002$). The positive trend
  indicates partial support for the structural account.}
\end{subfigure}
\hfill
\begin{subfigure}[t]{0.48\textwidth}
  \centering
  \includegraphics[width=\textwidth]{P1_03_interval_bars.png}
  \caption{Mean hue distance by interval type. The dashed line at
  180$^\circ$ marks the complementary-color prediction for the tritone
  and perfect fifth. Neither interval approaches this value.}
\end{subfigure}
\caption{Structural account: interval distance predicts hue distance,
but the strong complementary-hue prediction fails.}
\label{fig:structural}
\end{figure}

Among the six formalized structural theories, the two statistics
dissociate informatively. On internal geometry alignment
(Equation~\ref{eq:geometry}), four theories exceed the random-hue 95th
percentile of $\rho_{\text{geom}} = .23$: IRP and SSC realize a
perfect interval-to-arc isomorphism by construction
($\rho_{\text{geom}} = 1.00$, empirical $p < .001$), followed by MPA
($.68$, $p < .001$) and FRI ($.57$, $p < .001$); PDM ($.06$,
$p = .28$) and HEM ($-.12$, $p = .85$) fail. On angular error against
the pooled historical consensus, however, only IRP
(46.2$^\circ$, Monte Carlo $p = .001$) and FRI (54.9$^\circ$,
$p = .009$) score significantly below the 90$^\circ$ random-hue null;
HEM (70.7$^\circ$, $p = .10$), MPA (75.8$^\circ$), PDM
(76.3$^\circ$), and SSC (81.2$^\circ$) do not
(Figure~\ref{fig:theories}, left). The pattern is coherent: theories
that both encode the isomorphism \emph{and} anchor it to the spectrum
(IRP, FRI) match the historical record, whereas SSC shows that an
isomorphism anchored elsewhere lands far from the consensus, and
theories routing through intermediate perceptual dimensions (PDM,
HEM) capture neither. Notably, IRP and FRI---derived independently,
through different mathematical formalisms---cluster together in the
theory dendrogram (Figure~\ref{fig:theories}, right; pairwise distance
40.1$^\circ$). Two independent theories landing on the same
below-chance geometry is good evidence that a structural isomorphism
signal is present in the historical corpus, even if it is a partial
one.

\begin{figure}[htbp]
\centering
\begin{subfigure}[t]{0.52\textwidth}
  \centering
  \includegraphics[width=\textwidth]{P1_05_theory_errors.png}
  \caption{Consensus angular error for the six structural
  theories. Only IRP and FRI score significantly below the random-hue
  null ($\sim$90$^\circ$); HEM, MPA, PDM, and SSC do not. Vertical
  lines mark spectral/fifths benchmarks and the random-hue null.}
\end{subfigure}
\hfill
\begin{subfigure}[t]{0.45\textwidth}
  \centering
  \includegraphics[width=\textwidth]{P1_05_theory_dendrogram_beautified.png}
  \caption{Dendrogram of the six theories by mutual geometry distance.
  IRP and FRI form a tight cluster at low linkage distance, confirming
  their structural kinship.}
\end{subfigure}
\caption{Structural theory evaluation. Only direct isomorphism theories
(IRP, FRI) achieve below-chance geometry error.}
\label{fig:theories}
\end{figure}

The pairwise structure of these theories is shown directly in
Figure~\ref{fig:theory_heatmap}, and the predicted hue assignments of
each theory, overlaid on the observed global consensus, in
Figure~\ref{fig:theory_wheel}. IRP is the closest match to the
historical consensus among the six first-principles theories (the
spectral benchmark remains closer still), though it diverges
substantially in the middle register
(D--G$\sharp$) - a reminder that even the best structural theory
captures ordering without pinning down assignment.

\begin{figure}[htbp]
\centering
\includegraphics[width=0.72\textwidth]{P1_05_theory_heatmap.png}
\caption{Pairwise alignment heatmap for the six structural
theories. Cell color encodes the mean angular distance between the hue
assignments of each pair of theories. IRP and FRI form the tightest
pair (40.1$^\circ$), confirming that they capture the
same structural intuition through different mathematical formalisms.}
\label{fig:theory_heatmap}
\end{figure}

\begin{figure}[htbp]
\centering
\includegraphics[width=0.72\textwidth]{P1_05_theory_wheel.png}
\caption{Predicted hue assignments for the six structural theories
overlaid on the global consensus wheel. IRP is the closest match to
the observed historical consensus, though it still diverges
substantially in the middle register (D--G$\sharp$).}
\label{fig:theory_wheel}
\end{figure}

\subsection{Testing the Semantic Account: Is the Correspondence Carried
by Language?}

\textbf{Method.} We probed three pretrained sentence-transformer
models - MiniLM \citep{wang2020}, MPNet \citep{song2020}, and Nomic
Embed \citep{nussbaum2024} - to ask whether the linguistic descriptions of
pitches and colors carry pitch--hue structure on their own. Pitch
classes (C through B) were embedded using six descriptor levels: bare
name, musical context (scale degree, key relationships), perceptual
(frequency, register), timbral (instrument descriptions), connotative
(emotional/associative character), and solf\`{e}ge syllable. Two
independent descriptor sets were used: a primary \emph{LLM-generated}
set produced with an explicit constraint against
color, brightness, and temperature words (enforced by an automated
validator against a 35-term blocklist), so that any recovered geometry
could not be a trivial artifact of color words leaking into the pitch
descriptions, and a secondary \emph{factual} set derived
deterministically from music-theory constants (frequency, scale
degree, interval structure, instrument registers), used as a
robustness check on descriptor wording. Color-side embeddings used names from the XKCD Color
Survey \citep{munroe2010xkcd} - 954 crowdsourced color names with sRGB
hex codes. For each pitch class, the top-$k$ nearest XKCD colors (by
cosine similarity) formed a centroid hue angle, and the resulting
12-point pitch--hue mapping was scored by geometry alignment: the
Spearman $\rho$ between all $\binom{12}{2} = 66$ pairwise circular
pitch-step distances and the corresponding pairwise hue angular
distances.

\textbf{Results.} The sentence-transformer probing yields a clear and interpretable
result (Table~\ref{tab:lm_probing}). The strongest geometry signal is
Nomic with timbral descriptors ($\rho = .368$, $p = .007$):
statistically significant, well above the random-hue null, and
consistent across permutation iterations. So the semantic account does
capture something, but only partly---the timbral signal remains well
below the spectral geometry baseline ($\rho = 1.000$). One caveat
governs the interpretation throughout: timbral vocabulary is heavily
affective, so this signal cannot be read as purely linguistic. We treat
it as a semantics--affect entangled channel and return to the
distinction in the Discussion.

\begin{table}[htbp]
\centering
\caption{Sentence-transformer probing results. Geometry alignment is
Spearman $\rho$ between pairwise circular pitch-step distances and
pairwise hue angular distances; $p$-values are from 10,000
label-shuffle iterations, computed against condition-specific nulls
(the random-hue row is a global geometric reference, not the null used
for the $p$-values). The top conditions of the 3~(model)
$\times$~6~(descriptor) grid are shown; the full 18-cell grid, with
Benjamini--Hochberg corrected $q$-values, is given in Supplementary
Table~S4.}
\label{tab:lm_probing}
\begin{tabular}{@{}llcc@{}}
\toprule
\textbf{Model} & \textbf{Descriptor level} &
  \textbf{Geometry $\rho$} & \textbf{Perm.\ $p$} \\
\midrule
Nomic  & Timbral         & \textbf{0.368} & .007 \\
Nomic  & Perceptual      & 0.269          & .044 \\
MPNet  & Perceptual      & 0.203          & .048 \\
Nomic  & Bare names      & 0.173          & .070 \\
MPNet  & Musical context & 0.156          & .060 \\
MPNet  & Connotative     & 0.124          & .124 \\
MPNet  & Bare names      & 0.055          & .226 \\
MiniLM & Solf\`{e}ge     & 0.019          & .344 \\
\midrule
\multicolumn{2}{l}{\textit{Spectral geometry baseline}} & 1.000 & -- \\
\multicolumn{2}{l}{\textit{Circle-of-fifths baseline}} & $-0.191$ & -- \\
\multicolumn{2}{l}{\textit{Random-hue null (mean / 95th pct.)}}
  & $-0.001$ / 0.235 & -- \\
\bottomrule
\end{tabular}
\end{table}

The descriptor-level pattern points to the mechanism
(Figure~\ref{fig:semantic}, left). Timbral and perceptual descriptors
outperform bare names and solfège across all three models. Timbral
language---``warm'', ``round'', ``bright'', ``cutting''---carries
affective content that links pitch and color through shared experiential
vocabulary, and the contrast with the failed solfège condition (MiniLM:
$\rho = .019$, $p = .344$) makes pitch-name familiarity an unlikely
confound. The PCA projection (Figure~\ref{fig:semantic}, right) shows
how: Nomic's timbral pitch embeddings separate lower from higher pitches
along a dimension encoding ``warm/round'' versus ``bright/cutting'',
which is the geometric substrate of the Spearman alignment. In short,
the embedding space encodes that high notes are described as bright and
that bright things are described as blue.

\begin{figure}[htbp]
\centering
\begin{subfigure}[t]{0.54\textwidth}
  \centering
  \includegraphics[width=\textwidth]{P1_07v2_lm_error_bars.png}
  \caption{Geometry alignment by descriptor level and model
  (LLM-generated descriptors). Nomic shows the strongest pitch--hue
  geometry under timbral/perceptual descriptor conditions.}
\end{subfigure}
\hfill
\begin{subfigure}[t]{0.43\textwidth}
  \centering
  \includegraphics[width=\textwidth]{P1_07_latent_pitch.png}
  \caption{PCA projection of Nomic pitch-descriptor embeddings (timbral
  level). PC1 separates lower from higher pitches; PC2 encodes
  ``warm/round'' vs.\ ``bright/cutting'' timbral quality.}
\end{subfigure}

\vspace{0.75em}
\begin{subfigure}[t]{0.6\textwidth}
  \centering
  \includegraphics[width=\textwidth]{P1_07_latent_colors.png}
  \caption{Companion PCA projection of the 954 XKCD color-name
  embeddings (Nomic), each point colored by its actual hex value. Hue
  families form coherent clusters, showing that the color-side space
  against which pitch embeddings are matched itself encodes hue
  structure learned purely from text.}
\end{subfigure}
\caption{Semantic account: sentence-transformer probing reveals
significant pitch--hue geometry under timbral descriptor conditions,
consistent with a semantics--affect entangled channel.}
\label{fig:semantic}
\end{figure}

Results with the factual descriptor set are only partially consistent,
and we report them as a genuine constraint on interpretation. Several
factual conditions reach significance---MPNet with musical-context
descriptors is the strongest ($\rho = .334$, permutation $p = .012$),
and MiniLM ($\rho = .265$, $p = .048$) and MPNet ($\rho = .197$,
$p = .034$) recover geometry under timbral descriptors---but the
model--level profile shifts: Nomic's timbral signal, the strongest
cell of the primary grid, is absent under factual wording
($\rho = -.01$). The full factual grid is given in Supplementary
Table~S4. Geometry alignment is therefore recoverable from more than
one descriptor family, but which model--level cell carries it depends
on wording---a sensitivity consistent with the entangled
semantics--affect reading rather than with a wording-independent
property of the embedding spaces.
The per-model predicted mappings are shown in
Figure~\ref{fig:lm_wheels}. Figure~\ref{fig:crossling} reports two
controls: the cross-linguistic panel is a robustness check showing that
pitch--hue geometry remains weak across naming conventions, arguing
against name-specific artifacts; and the emotion-mediation panel's
small positive correlation ($\rho = .139$, $p < .001$) is consistent
with the hedonic mediation examined next.

\begin{figure}[htbp]
\centering
\includegraphics[width=\textwidth]{P1_07_lm_wheels.png}
\caption{Predicted hue assignments from the three language models
(MiniLM, MPNet, Nomic) under LLM-generated timbral descriptors. Nomic
shows the most structured prediction (warm reds at lower pitches, cooler
hues at higher pitches), consistent with its higher Spearman geometry
alignment ($\rho = .368$, $p = .007$). MiniLM predictions are noisier,
reflecting its smaller training corpus.}
\label{fig:lm_wheels}
\end{figure}

\begin{figure}[htbp]
\centering
\includegraphics[width=\textwidth]{P1_07v2_crossling_emotion.png}
\caption{Cross-linguistic naming and emotion mediation controls. Left:
geometry alignment for five naming conventions (English, Japanese,
German, solf\`{e}ge with context, bare solf\`{e}ge), a robustness check
showing that pitch--hue geometry is weak across all naming conventions.
Right: emotion mediation test - the small positive Spearman correlation
($\rho = .139$, $p < .001$) indicates that embedding proximity between
pitch and color terms tracks shared emotional valence, consistent with
the hedonic account.}
\label{fig:crossling}
\end{figure}

\subsection{Corpus-Level Regularities: Testing the Statistical Account
and Its Neighbors}

\textbf{Method.} Ten corpus-level hypotheses were evaluated on the
pooled consensus and per-pitch statistics of the 25 retained systems
(Table~\ref{tab:statistical_h}). Four operationalize the environmental
co-occurrence account proper ($H_{b4}$, $H_{b6}$, $H_{b7}$, $H_{b9}$);
the remainder test neighboring accounts and are attributed accordingly
in the table: $H_{b1}$, $H_{b3}$, and $H_{b8}$ probe structural
predictions, $H_{b2}$ probes a salience/affect asymmetry, $H_{b5}$
and $H_{b10}$ probe notation-driven and tradition-driven cultural
variation. The Rayleigh
test \citep{mardia1972} was used for the non-uniformity prediction
($H_{b4}$); because $H_{b4}$ restates the non-uniformity already
reported in Section~\ref{sec:global}, it is counted once in the
overall tally. Correlational hypotheses use Pearson tests over the 12
pitch classes or the interval table; group-difference hypotheses use
Welch $t$-tests.

\textbf{Results.} The battery produces three passes, six failures, and
one inconclusive result (Table~\ref{tab:statistical_h}). The clearest result is $H_{b6}$: the
correlation between interval size and between-system variance is
$r = -.94$, $p = .001$---though, because circular pitch-class distance
takes only six distinct values, this correlation is computed over
seven interval rows and its precision should be read accordingly.
Historical
mappers agree more about large intervals and disagree more about
semitones, the pattern one would expect from statistical learning of a
noisy environmental signal that constrains global ordering but not
fine-grained assignment. $H_{b4}$ adds support: hue assignments
cluster toward warm colors, consistent with the environmental
co-occurrence of low-frequency sounds and warm atmospheric conditions.
$H_{b7}$, by contrast, is downgraded from the ``pass'' a naive reading
would give it: the fifth--tritone difference is 5.0$^\circ$ against
standard deviations of 22--26$^\circ$, was not formally tested, and the
fifth is circularly equivalent to the fourth; we report it as
directionally consistent but inconclusive.

\begin{table}[htbp]
\centering
\caption{Corpus-level hypothesis test results. The Account column
attributes each hypothesis to the theoretical account it
operationalizes; only the Statistical rows bear directly on the
environmental co-occurrence account.}
\label{tab:statistical_h}
\begin{tabular}{@{}p{0.9cm}p{4.6cm}p{1.7cm}p{1.9cm}p{3.2cm}@{}}
\toprule
\textbf{Hyp.} & \textbf{Description} & \textbf{Account} &
  \textbf{Result} & \textbf{Key statistic} \\
\midrule
$H_{b1}$ & Lightness increases with chromatic pitch order
  & Structural & Fail & $r = -.53$, $p = .076$ \\
$H_{b2}$ & Agreement is higher for high-chroma assignments
  & Hedonic & Fail & $r = .55$, $p = .063$ \\
$H_{b3}$ & Consonant intervals show larger hue distance than dissonant
  & Structural & Fail & $t = 1.05$, $p = .34$ \\
$H_{b4}$ & Hue distribution non-uniform
  & Statistical & \textbf{Pass} & Rayleigh $R = 0.66$, $p = .003$ \\
$H_{b5}$ & Sharps differ from naturals in lightness or hue
  & Cultural & Fail & $p = .53$ / $.95$ \\
$H_{b6}$ & Interval size predicts variance
  & Statistical & \textbf{Pass} & $r = -.94$, $p = .001$ \\
$H_{b7}$ & Perfect fifth $<$ tritone hue distance
  & Statistical & Inconclusive & $80.7^\circ < 85.7^\circ$ (untested) \\
$H_{b8}$ & Primary triad (C--E--G) hues $\sim$120$^\circ$ apart
  & Structural & Fail & mean $66.4^\circ$, dev.\ $53.6^\circ$ \\
$H_{b9}$ & Agreement correlates with consistency
  & Statistical & \textbf{Pass} & $r = -.62$, $p = .03$ \\
$H_{b10}$ & Tradition-specific fifths differ from pooled
  & Cultural & Fail & $t = 0.15$, $p = .89$ \\
\bottomrule
\end{tabular}
\end{table}

Two failures deserve comment. The failure of $H_{b1}$ is the more
surprising: pitch--lightness is among the most robust cross-modal
correspondences in the behavioral literature
\citep{marks1978,ward2006}, yet across the historical consensus the
correlation between chromatic pitch order and assigned lightness is
directionally \emph{negative} ($r = -.53$, $p = .076$). The historical
record encodes hue structure, not the pitch--lightness gradient---%
consistent with our inclusion criteria, which excluded
lightness-only systems, but also a caution against reading historical
proposals as transcriptions of perceptual correspondence. The failure
of $H_{b10}$ ($p = .89$) shows that the perfect-fifth hue distance is
statistically indistinguishable across reasoning traditions; combined
with the family analysis, it suggests that traditions differ in
\emph{which} colors they assign more than in how they space
harmonically related pitches.

\subsection{Testing the Hedonic Account: Do Pitch and Color Share an
Affective Geometry?}

\textbf{Method.} The hedonic account predicts that the organization of
auditory and visual dimensions in a general sensorimotor
representation is captured by affective structure. To test this, the
cross-modal dimensions in the CMC database of the following section
were each coded on 11 sensorimotor--affective dimensions (visual,
auditory, tactile, gustatory, olfactory, interoceptive, motor,
valence, arousal, magnitude, spatial), an inventory modeled on the
Lancaster Sensorimotor Norms \citep{lynott2020}; because the Lancaster
norms cover words rather than the technical dimension labels used
here, the profiles were coded by the authors, a construct limitation
weighed in the Discussion. The dimension profiles were reduced to
three principal components (68.3\% of variance). Ten dimensions
central to the pitch--hue question---pitch height, loudness, tempo,
hardness on the auditory--tactile side; lightness, brightness,
saturation, hue warmth on the visual side; size and pleasantness as
cross-modal anchors---were then assigned coordinates in
Evaluation-Potency-Activity (EPA) affective space \citep{osgood1957},
drawing on published semantic-differential and color-affect norms
\citep{ou2004,hevner1935,wedin1972}. Procrustes
alignment \citep{gower1975} found the rotation, reflection, and uniform
scaling minimizing squared distances between the sensorimotor PCA
configuration and the EPA configuration of the same ten dimensions.
The \textbf{Procrustes disparity} is
\begin{equation}
  d^2 = 1 - \left(\sum_{k=1}^{p} \sigma_k\right)^2
\label{eq:procrustes}
\end{equation}
where $\sigma_1, \ldots, \sigma_p$ are the singular values of
$\mathbf{X}^\top \mathbf{Y}$, with $\mathbf{X}$ and $\mathbf{Y}$ the
centered, unit-scaled configuration matrices; a disparity
of $0$ denotes perfect alignment and $1$ no alignment. Significance was
assessed by permutation of the correspondence assignments before
alignment (10,000 permutations); the exact $p$-value and the mean and
standard deviation of the null disparity distribution are reported in
Supplementary Table~S5, along with dimension-wise alignment statistics.

\textbf{Results.} The hedonic account gives the best result of the four. The
Procrustes disparity between the sensorimotor and EPA configurations is
0.559, significantly below the random-correspondence permutation null
(null mean $0.833$, SD $0.106$; $p = .027$) and the best
single-account result across the four
theoretical tests, indicating that the space in which auditory and
visual dimensions are organized shares an affective
geometry. Dimension-wise correlations between matched coordinates
after alignment identify Evaluation as the dominant channel
($r = .74$, $p = .014$), with Potency ($r = .47$, $p = .17$) and
Activity ($r = .37$, $p = .29$) weaker and individually
non-significant (Supplementary
Table~S5). Figure~\ref{fig:hedonic} shows the configuration.

\begin{figure}[htbp]
\centering
\includegraphics[width=\textwidth]{P1_06_walker_epa.png}
\caption{EPA Procrustes analysis. Left: the ten cross-modal
dimensions positioned in affective space (Evaluation vs.\ Activity
shown), with auditory dimensions (blue), visual dimensions (orange),
and cross-modal anchors (purple); the auditory and visual clusters
overlap along the Evaluation axis. Middle: the Procrustes disparity
of 0.559 sits between perfect alignment (0) and the permutation null
($0.833$); the pitch--brightness correspondence, shown for
calibration, aligns more closely than pitch--hue. Right: mean
behavioral effect sizes of the catalogued CMC pairings grouped by
proposed mechanism; pitch--hue ($d = 0.20$) is the weakest documented
pairing.}
\label{fig:hedonic}
\end{figure}

Evaluation ($E$) is the dominant aligning dimension: the good--bad,
pleasant--unpleasant axis orders auditory and visual dimensions in the
same way, so that what aligns pitch-space with color-space most
reliably is shared hedonic value. Activity and Potency contribute
positively but weakly. Converging evidence at the item level comes
from the embedding analysis of the previous section: proximity between
pitch and color terms in embedding space tracks shared
valence--arousal profiles ($\rho = .139$, $p < .001$;
Figure~\ref{fig:crossling}, right). This pattern is consistent with
Palmer and colleagues' \citep{palmer2013} demonstration of
emotion-mediated music--color associations, and locates the mediation
on the Evaluation axis: what seems to matter is the shared
pleasantness of sounds and colors more than their shared energy or
power.

\subsection{Positioning Pitch--Hue in the CMC Landscape}

\textbf{Method.} A structured database of 34 documented CMC pairings
was compiled from the published literature (sources in Supplementary
Table~S6), each pairing coded for source and target modality,
direction, mechanism proposed in the source literature, and a
literature-derived behavioral effect size on a common Cohen's-$d$
scale with a plausible range. Each participating dimension was coded
on the 11 sensorimotor--affective dimensions described in the previous
section; cosine distance between these
profiles quantifies dimensional remoteness
between paired modalities, and the landscape was visualized with UMAP
\citep{mcinnes2018}. Because both the dimension profiles and the
effect sizes are coded from the literature rather than measured, this
analysis is a quantitative synthesis of published claims, not new
measurement.

\textbf{Results.} This analysis contextualizes the preceding ones.
Sensorimotor
distance correlates negatively with behavioral effect size across the
34 pairings ($r = -.379$, $p = .027$): the further apart two modalities
are in sensorimotor space, the weaker the correspondence between them.
Pitch and hue have one of the largest pairwise sensorimotor distances
in the auditory--visual CMC literature, and pitch--hue occupies an
isolated peripheral position in the UMAP embedding. Its weakness
relative to pitch--elevation or loudness--size therefore carries
theoretical information rather than being an accident: the modalities
share no direct prothetic dimension, so the correspondence must be
brokered by higher-order structure.

% ============================================================
\section{Discussion}

\subsection{Is There Structure in Three Centuries of Disagreement?}

Yes, but the structure is family-specific and affectively mediated
rather than universal or physically determined. Global agreement is low
(mean 0.22), confirming that no pitch--hue mapping commands
cross-cultural authority. At the same time, hue assignments are
non-uniformly distributed (Rayleigh $R = 0.66$, $p = .003$), the
Pedagogical and Synesthetic families sit above the flat baseline while
the Cosmological family sits below it, and this ordering---though not
individually significant at current family sizes---is coherent across
every way we sliced the corpus.

The central methodological contribution of this paper is establishing
that the reasoning family, not the individual system, is the correct
unit of analysis for the historical corpus. All prior studies computing
``agreement'' have pooled heterogeneous traditions and observed global
disagreement, thereby missing the structure hiding within each tradition.
The eight-family classification dissolves the apparent contradiction
between ``correspondences exist'' (within families, they do) and ``no
consensus exists'' (across families, there isn't one). This is
meta-analytic heterogeneity in its
classic form: pooling across subgroups with different underlying
mappings masks the agreement present within each subgroup. The
Cosmological family sharpens the point from the other direction---its
within-tradition agreement falls \emph{below} the pooled baseline, so
grouping by tradition does not merely recover hidden structure but
also exposes traditions that carry none. One honesty constraint
applies throughout: with 3--5 systems in most families, neither the
individual family deviations nor the global heterogeneity statistic
survives permutation testing (Supplementary Table~S2), so the family
ordering is an effect-size-level pattern awaiting confirmation on an
expanded corpus, not an established difference.

\subsection{Which Account Best Explains the Structure?}

Table~\ref{tab:account_summary} collects the verdicts.

\begin{table}[htbp]
\centering
\caption{Summary of the four theoretical accounts against the
evidence.}
\label{tab:account_summary}
\begin{tabular}{@{}p{2.4cm}p{4.4cm}p{3.6cm}p{3.4cm}@{}}
\toprule
\textbf{Account} & \textbf{Key prediction} & \textbf{Key statistic} &
  \textbf{Verdict} \\
\midrule
Structural & Interval geometry maps to hue geometry
  & $r = .393$; IRP 46.2$^\circ$ vs.\ 90$^\circ$ null; fifths at
    80.7$^\circ$, not 180$^\circ$
  & Partial: ordering, not assignment \\
Semantic & Shared language carries the mapping
  & Timbral $\rho = .368$, $p = .007$; solf\`{e}ge $\rho = .019$
  & Contributing, entangled with affect \\
Statistical & Co-occurrence constrains color spacing
  & $r = -.94$ (interval size vs.\ variance); Rayleigh $R = 0.66$
  & Supported for large-scale structure only \\
Hedonic & Shared affect aligns the modalities
  & Procrustes disparity 0.559, $p = .027$; Evaluation dominant
  & Best supported \\
\bottomrule
\end{tabular}
\end{table}

\textbf{Structural account: partial support.} The semitone--hue
distance correlation ($r = .393$) and the below-chance consensus error
of IRP and FRI establish that interval structure leaves a detectable
imprint on historical color assignments. But the flagship prediction
fails: perfect fifths separate by 80.7$^\circ$, not 180$^\circ$, and
the interval--distance function saturates well before reaching the
complementary region. Structural regularities constrain the ordinal
relationship between interval size and hue distance---a partial
scaffold---but cannot fix specific pitch-to-hue assignments. Only the
spectrally anchored isomorphism theories (IRP, FRI) match the
historical consensus; an isomorphism anchored elsewhere (SSC) or
routed through intermediate perceptual dimensions (MPA, PDM, HEM)
does not. The
structural signal in the data is an isomorphism signal: narrow, but
genuine.

\textbf{Semantic account: contributing but insufficient.} Language
models encode significant pitch--hue geometry under timbral descriptor
conditions ($\rho = .368$, $p = .007$). The active channel is shared
\emph{experiential} vocabulary---``warm'', ``bright'', ``cutting''---rather
than arbitrary name co-occurrence, since solfège labels produce nothing
($\rho = .019$). The effect is genuine but incomplete: timbral language
carries heavy affective content, so the geometry signal under timbral
conditions is partly hedonic. The semantic and hedonic accounts cannot
be fully separated here, which is itself informative---language and
affect are closely tied in the pitch--hue domain.

\textbf{Statistical account: specific support for large-scale
structure.} The $r = -.94$ correlation between interval size and
between-system variance is the most tightly predicted result in the
analysis, albeit one computed over the six distinct circular interval
distances. Environmental statistics constrain global ordering (distant pitches get
distant colors) but not fine-grained assignment. The account's
explanatory scope is thus specific: it tells you how far apart the
colors should be, but not what they should be.

\textbf{Hedonic account: the best-supported.} The EPA Procrustes
disparity (0.559, $p = .027$) is the strongest single-account result
across the four
tests. Evaluation is the primary aligning dimension ($r = .74$), so
auditory and visual dimensions
are connected most reliably through their shared hedonic value. The
item-level embedding evidence points the same way: pitch--color term
proximity tracks shared valence--arousal profiles ($\rho = .139$). The
family
data fit this reading: Pedagogical systems converge because educators
independently settle on what \emph{feels} right to learners, without
sharing a physical theory, and Synesthetic correspondences appear to tap
the same affective architecture as non-synesthetic CMCs---an
amplification of a signal that is already present, rather than a
separate system. Palmer and colleagues \citep{palmer2013} showed that
music--color associations are emotion-mediated; our results locate
that mediation on the Evaluation (good--bad) axis, with Potency and
Activity contributing only weakly.

Pulling these together, pitch--hue correspondences appear to rest on a
\emph{shared affective geometry}, partly scaffolded by structural
regularities and environmental statistics, faintly echoed in language
under timbral conditions, and not recoverable from any single
first-principles map. The hedonic account is necessary but not
sufficient: without the structural scaffolding, the data's
interval--distance relationship has no explanation.

\subsection{Positioning Pitch--Hue Among Cross-Modal Correspondences}

Pitch--hue is among the hardest auditory--visual CMCs to explain, and
the sensorimotor landscape analysis suggests why. Pitch and hue are
dimensionally remote: they share no prothetic (intensive) dimension of
the kind that directly grounds pitch--size, pitch--brightness, and
loudness--size correspondences \citep{marks1978,walsh2003}. The
$r = -0.379$ correlation between sensorimotor distance and behavioral
effect size ($p = .027$) captures this: the further apart two modalities
are in sensorimotor space, the weaker their correspondence. Pitch--hue
sits at the far end of this trade-off, so any correspondence between
them has to be brokered by higher-order structure---affective valence,
cultural co-occurrence, symbolic convention---which is what the present
results show. The peripheral UMAP position of pitch--hue in the
auditory--visual cluster then follows from the sensorimotor distance
account rather than being an anomaly.

\subsection{Implications for Neural Congruency Research}

The historical analysis exposes a critical unresolved problem in the
prior literature: no principled, data-driven criterion for pitch--hue
congruency exists that is independent of behavioral report. The field
has been defining congruency by asking participants what feels
congruent - and no neural correlate of pitch--hue congruency has been
identified as a result. The present findings break this circularity
directly. Because no single historical map commands ground-truth
authority, congruency definitions for neuroimaging experiments must not
be derived from specific historical systems such as Newton's or
Scriabin's. Doing so imports an arbitrary cultural choice into the
experimental design and biases RSA matrices against the actual
structure in the data.

The right approach is experimentally defined congruency - for example,
audiovisual Stroop paradigms in which participants hear a tone while
viewing a written note name that either matches or conflicts with the
sounded pitch. The EPA results support this choice precisely: the
Stroop conflict manipulation creates a mismatch on the Evaluation
dimension that this analysis identifies as the primary
affective channel of the pitch--hue correspondence. Beyond the Stroop,
the theoretical RDMs derived from IRP, FRI, and EPA alignment can be
submitted directly to RSA against EEG or fMRI embeddings. For the
first time, a neural investigation of pitch--hue processing has a
quantitatively grounded, theory-derived model to test.

\subsection{Limitations and Future Directions}

Several limitations constrain the current conclusions, and we state
them plainly because they bound what the analysis can and cannot
support. First, and most important, the family-level results rest on
small samples. Several analyzed families contain only a handful of
systems; each family score carries a
bootstrap interval and a label-permutation test (Supplementary
Table~S2), and none of the family deviations---nor the global
heterogeneity statistic ($p = .81$)---reaches significance under
permutation. The apparent ordering of
families is therefore suggestive rather than established, and
replication on an expanded catalogue is needed before the family
ranking can be treated as settled. Second, the study tests a
large battery of hypotheses - six interval hypotheses, six structural
theories, ten corpus-level predictions, and multiple model--descriptor
combinations - so individually significant results must be read in the
context of this wide inferential surface. We designated four primary
claims in advance of the confirmatory write-up (family heterogeneity,
the EPA Procrustes alignment, the interval--variance correlation, and
the timbral embedding alignment) and collect all primary tests with
Benjamini--Hochberg FDR-corrected $q$-values in Supplementary
Table~S7. Of the four primary claims, three survive correction at
$q < .05$ (the EPA alignment, $q = .045$; the interval--variance
correlation, $q = .007$; the timbral embedding alignment, $q = .024$),
as do the semitone--hue correlation, the Rayleigh non-uniformity, the
agreement--consistency correlation, and the effect-size--distance
relationship; the family-heterogeneity permutation tests do not
($q \ge .13$).
A pre-registered confirmatory study would nonetheless put the
surviving effects on firmer footing. Third, color
standardization introduces measurement error that is difficult to
quantify: resolving free-text and pigment-name specifications to CIELAB
via crowd-sourced and reference databases is principled but noisy, and
some of the low global agreement may reflect conversion noise rather
than genuine disagreement. Fourth, the hedonic test operates on
author-coded constructs at both ends: the 11-dimensional
sensorimotor profiles of the cross-modal dimensions were coded by the
authors on an inventory modeled on the Lancaster norms rather than
taken from participant ratings, and the EPA coordinates of the ten
aligned dimensions were assembled from heterogeneous published norms.
The Procrustes result is significant under permutation, but both
configurations are inferential constructs over a small term set
($n = 10$), and the strength of the result should be weighed against
this construct uncertainty; replacing the coded profiles with
participant-normed ratings is a direct next step.
Fifth, the catalogue is not a random sample:
systems that were never published, or specified colors in forms that
could not be standardized, are excluded, and the retained systems skew
toward formally specified, Western-centric proposals; the family
agreement scores are therefore lower bounds on within-family coherence
for well-documented traditions, not estimates of a universal population
parameter. The EPA norms are likewise predominantly Western, so the
Evaluation-first finding should be tested cross-culturally before
being generalized. Finally, the sentence-transformer probing demonstrates
geometry alignment rather than specific mapping prediction, and a
positive timbral result does not rule out that the signal is driven by
affective content in the descriptors rather than purely linguistic
co-occurrence; the hedonic and semantic accounts overlap here, and
disentangling them requires targeted descriptor construction.

These limitations also locate the contribution. This is a meta-analytic
and methodological study - a curated catalogue, a reference-free
agreement framework, and a direct comparison of accounts---rather
than a behavioral experiment, and it should be read as one. Its purpose
is to extract the structure latent in the historical record and to hand
the field a quantitatively grounded, theory-derived target for the
experiments that must follow. The three most important next steps are:
(1)~extend the catalogue through non-Western archival sources and
non-published practitioner traditions, and report interval estimates on
the expanded sample, to reduce both sampling bias and statistical
fragility; (2)~collect large-scale behavioral pitch--hue matching data
across multiple cultural backgrounds to test whether the family-level
historical structure predicts individual differences; and (3)~deploy
the IRP, FRI, and EPA-aligned theoretical RDMs as predictors in RSA of
EEG and fMRI responses, testing directly whether neural geometry
mirrors the structure identified here.

\subsection{Conclusion}

Three centuries of pitch--hue proposals do not converge on a universal
mapping, but they are far from random. Organized by reasoning tradition,
they show a within-family ordering that prior flat analyses miss:
Pedagogical systems reach agreement 57\% above the global baseline,
Synesthetic systems converge through a constrained affective
architecture, and Cosmological systems, despite their ambition, agree
less among themselves than the pooled corpus does---an ordering that
awaits confirmation on a larger corpus. Of the four theoretical
accounts, the
hedonic one is best supported---auditory and visual dimensions share
an affective geometry, with Evaluation, the good--bad axis, as the
primary bridge---while structural
regularities provide a partial scaffold, environmental statistics
constrain global ordering, and language carries a secondary timbral
signal. None of these accounts is sufficient on its own; together they
offer a plausible account of why proposals from Pythagorean cosmologists
to Suzuki educators land where they do. The weakness of pitch--hue
relative to other auditory--visual CMCs follows from the sensorimotor
distance account rather than being a disappointment: pitch and hue are
far apart in the sensorimotor world, so any correspondence between them
has to be carried by affective structure---and on the present evidence,
shared hedonic value is what does most of that work.

% ============================================================
\section*{Supplementary Material}

The Supplementary Material (\texttt{supplementary\_material.pdf})
comprises: \textbf{S1}, the full 70-system
catalogue (system, year, reasoning family, retention status,
exclusion reason); \textbf{S2}, family agreement
bootstrap intervals, label-permutation $p$-values, the global
heterogeneity test, and the simulated
chance distribution of the agreement score; \textbf{S3}, the comparison
of the Pedagogical (F8) consensus mapping against canonical
spectral order; \textbf{S4}, the full 3~$\times$~6 sentence-transformer
probing grid for both descriptor sets, with FDR-corrected $q$-values;
\textbf{S5}, exact Procrustes permutation statistics and dimension-wise
alignment; \textbf{S6}, the coded CMC database with effect sizes and
key sources for the 34
pairings; and \textbf{S7}, the consolidated multiplicity table of
all primary tests.

\section*{Open Practices Statement}

The historical catalogue, CIELAB conversion pipeline, distance
matrices, family classification codes, and analysis scripts are
available at \texttt{[repository link]}.
% TODO(authors): create the public repository (e.g., OSF or GitHub)
% and insert the link; APP expects availability at submission.
None of the analyses were preregistered.

\section*{Conflict of Interest}

The authors declare no conflicts of interest.

% ============================================================
\bibliographystyle{apalike}
\bibliography{references}

\end{document}
